\begin{frame}
\frametitle{What is Parametricity}
\begin{block}{Danielsson, Hughes, Jansson \& Gibbons \cite{danielsson2006fast} tell us:}
\begin{quotation}
Functional programmers often reason about programs as if
they were written in a total language, expecting the results
to carry over to non-total (partial) languages. We justify
such reasoning.
\end{quotation}
\end{block}
\end{frame}

\begin{frame}
\frametitle{What is Parametricity}
\begin{block}{Philip Wadler \cite{wadler1989theorems} tells us:}
\begin{quotation}
Write down the definition of a polymorphic function on a piece of paper. Tell me its type, but be careful not to let me see the function's definition. I will tell you a theorem that the function satisfies.

The purpose of this paper is to explain the trick.
\end{quotation}
\end{block}
\end{frame}

\begin{frame}[fragile]
\frametitle{Types}
\begin{block}{first let's talk about types}
Suppose we encountered the following function definition:
\begin{lstlisting}
int add12(int)
\end{lstlisting}
\end{block}
\begin{itemize}
  \item<1> by the type alone, there are {$({2^{32}})^{2^{32}}$} possible implementations
  \item<2> but this is a significantly smaller number than \rotatebox{90}{8}
\end{itemize}
\end{frame}

\begin{frame}[fragile]
\frametitle{Types}
\framesubtitle{reading the code}
We might form a suspicion that \lstinline[style=scala]$add12$ adds twelve to its argument
\begin{lstlisting}[style=scala]
int `add12`(int)
\end{lstlisting}
\begin{tikzpicture}[remember picture,overlay]
\coordinate (aa) at ($(a1)+(7,2.0)$);
\node[note,draw,callout relative pointer={($(aa)-(11.2,-3.7)$)},right] at (aa) {\includegraphics[width=0.2\textwidth]{image/suspicion.jpg}};
\end{tikzpicture}
\end{frame}

\begin{frame}[fragile]
\frametitle{Types}
So we write some tests:
\begin{lstlisting}
add12(0)        = 12
add12(5)        = 17
add12(-5)       = 7
add12(223)      = 235
add12(5096)     = 5104
add12(2914578)  = 29145590
add12(-2914578) = -29145566
\end{lstlisting}
And conclude, yes, this function adds twelve to its argument
\end{frame}

\begin{frame}[fragile]
\frametitle{Types}
\framesubtitle{Woops!}
\begin{lstlisting}
def add12(n: Int): Int =
  if(n < 8000000) n + 12
  else n * 7
\end{lstlisting}
Narrowing the potential propositions about what this function does not do.
\end{frame}

\begin{frame}[fragile]
\frametitle{Types}
\begin{block}{another monomorphic example}
\begin{lstlisting}
List<int> function(List<int>)
\end{lstlisting}
\end{block}
\begin{itemize}
  \item<1> adds 17 to every 11th element?
  \item<2> drops every prime number?
\end{itemize}
\end{frame}

\begin{frame}[fragile]
\frametitle{Parametricity}
\begin{block}{a polymorphic example}
\begin{lstlisting}
<A> List<A> function(List<A>)
\end{lstlisting}
\end{block}
\begin{itemize}
  \item<1> this function returns elements in a list that always appear in the argument
  \item<2> this function never inspects the elements, only rearranges them
\end{itemize}
\end{frame}

\begin{frame}[fragile]
\frametitle{Parametricity}
\begin{block}{the goal}
\begin{itemize}
  \item<1> a significant number of possible things that this function does are eliminated, by no expenditure of effort
  \item<2> theorems about this function can be reliably constructed
\end{itemize}
\end{block}
\end{frame}

\begin{frame}[fragile]
\frametitle{Reasoning with parametricity}
\begin{block}{Fast and loose reasoning is morally correct \cite{danielsson2006fast}}
\begin{quotation}
Functional programmers often reason about programs as if
they were written in a total language, expecting the results
to carry over to non-total (partial) languages. We justify
such reasoning.
\end{quotation}
\end{block}
but what does this mean exactly?
\end{frame}

